% ============================================================================
% capibara_kawaii.tex  —  Capibara KAWAII / chibi ORIGINAL (obra propia, CC0)
% ----------------------------------------------------------------------------
% Personaje para la serie de matematica financiera (interes simple / compuesto).
% Estilo: chibi MUY tierno (tipo canal "EconoMonos"): plano, limpio y dulce.
% Cuerpo redondo/blandito marron claro pastel, orejitas pequenas redondas, ojos
% negros GRANDES con doble brillo, hocico crema con NARIZ grande redondeada
% (rasgo iconico del capibara), cachetes rosa suave y sonrisa feliz.
% Inspirado solo en el ESTILO kawaii (no es copia de ningun sticker con copyright).
%
% Requiere:  \usepackage{tikz}   (pgfkeys viene con tikz).  NO usa texto/fuentes.
%
% USO
%   \capibarakawaii                 % pose neutra, con rubor
%   \capibarakawaii[hold]           % bracitos arriba para sostener algo (naranja)
%   \capibarakawaii[blush=false]    % sin rubor
%   \capibarakawaii[wink]           % ojo izq cerrado (guino) -> para parpadeo
%
% RECOLOREAR (antes de llamar la macro, o editando los \providecolor de abajo):
%   \colorlet{capicuerpo}{brown!35}\colorlet{capioscuro}{brown!55}
%
% PARTES SEPARABLES (submobjects para animar en Manim / SVGMobject). Orden real
% de dibujo (pose neutra) -> conteo y mapa exactos en
% assets/svg-src/interes/CREDITOS.md. Parpadeo = escalar en y los ojos+brillos.
% ============================================================================

\providecolor{capicontorno}{HTML}{6E4E34}  % linea (marron calido suave)
\providecolor{capicuerpo}{HTML}{D9B38A}    % cuerpo (marron claro pastel)
\providecolor{capioscuro}{HTML}{C69A6A}    % orejas int / patitas (marron medio suave)
\providecolor{capihocico}{HTML}{F0DFC2}    % parche del hocico (crema suave)
\providecolor{capinariz}{HTML}{7A5638}     % nariz (marron calido)
\providecolor{capiojo}{HTML}{4A3526}       % ojos (marron muy oscuro, no negro duro)
\providecolor{capibrillo}{HTML}{FFFFFF}    % brillo de los ojos
\providecolor{capirubor}{HTML}{F3B7B0}     % rubor rosado suave

\tikzset{
  capilinea/.style={draw=capicontorno, line width=2.1pt,
                    line cap=round, line join=round},
}

\newif\ifcapiblush \capiblushtrue
\newif\ifcapihold  \capiholdfalse
\newif\ifcapiwink  \capiwinkfalse

\pgfkeys{
  /capi/.is family, /capi,
  blush/.is if=capiblush,
  hold/.is if=capihold,
  wink/.is if=capiwink,
}

\newcommand{\capibarakawaii}[1][]{%
  \begingroup
  \pgfkeys{/capi, blush=true, hold=false, wink=false, #1}%
  \begin{scope}
    % --- orejas pequenas y redondas (detras del cuerpo) ---
    \path[capilinea, fill=capicuerpo] (-0.92,1.30) circle[radius=0.34];
    \path[capilinea, fill=capicuerpo] ( 0.92,1.30) circle[radius=0.34];
    \path[fill=capioscuro] (-0.92,1.31) circle[radius=0.17];
    \path[fill=capioscuro] ( 0.92,1.31) circle[radius=0.17];
    % --- cuerpo (bolita redonda y blandita, marron claro pastel) ---
    \path[capilinea, fill=capicuerpo, rounded corners=1.35cm]
        (-1.50,-1.62) rectangle (1.50,1.42);
    % --- patitas cortas y redonditas (delante, abajo) ---
    \path[capilinea, fill=capioscuro]
        (-0.60,-1.60) ellipse[x radius=0.40, y radius=0.32];
    \path[capilinea, fill=capioscuro]
        ( 0.60,-1.60) ellipse[x radius=0.40, y radius=0.32];
    % --- bracitos cortos (delante) ---
    \ifcapihold
      \path[capilinea, fill=capicuerpo]
          (-0.68,-0.30) ellipse[x radius=0.28, y radius=0.42];
      \path[capilinea, fill=capicuerpo]
          ( 0.68,-0.30) ellipse[x radius=0.28, y radius=0.42];
    \else
      \path[capilinea, fill=capicuerpo]
          (-1.34,-0.66) ellipse[x radius=0.30, y radius=0.40];
      \path[capilinea, fill=capicuerpo]
          ( 1.34,-0.66) ellipse[x radius=0.30, y radius=0.40];
    \fi
    % --- cachetes rosa suave ---
    \ifcapiblush
      \path[fill=capirubor] (-1.00,-0.02) ellipse[x radius=0.26, y radius=0.17];
      \path[fill=capirubor] ( 1.00,-0.02) ellipse[x radius=0.26, y radius=0.17];
    \fi
    % --- ojos GRANDES con doble brillo (dulces) ---
    \ifcapiwink
      \draw[capilinea, line width=2.6pt]
          (-0.76,0.52) .. controls (-0.54,0.36) .. (-0.32,0.52);
    \else
      \path[fill=capiojo] (-0.54,0.50) circle[radius=0.35];
      \path[fill=capibrillo] (-0.43,0.66) circle[radius=0.135];
      \path[fill=capibrillo] (-0.66,0.40) circle[radius=0.070];
    \fi
    \path[fill=capiojo] ( 0.54,0.50) circle[radius=0.35];
    \path[fill=capibrillo] ( 0.65,0.66) circle[radius=0.135];
    \path[fill=capibrillo] ( 0.42,0.40) circle[radius=0.070];
    % --- hocico (parche crema suave) ---
    \path[fill=capihocico, rounded corners=0.45cm]
        (-0.52,-0.60) rectangle (0.52,0.08);
    % --- nariz grande redondeada (rasgo iconico del capibara, tierna) ---
    \path[capilinea, fill=capinariz, rounded corners=0.18cm]
        (-0.31,-0.16) rectangle (0.31,0.14);
    % --- boca: sonrisa feliz (dulce, curva hacia arriba) ---
    \draw[capilinea, line width=1.8pt]
        (0,-0.16) -- (0,-0.24)
        (-0.26,-0.24) .. controls (-0.10,-0.50) and (0.10,-0.50) .. (0.26,-0.24);
  \end{scope}
  \endgroup
}
